\documentclass{article}

%A mettre tout le temps
\usepackage{lcpackage}

\begin{document}
\pagestyle{fancy}
\fancyhf{}
\fancyhead[L]{{\fontfamily{phv}\selectfont Maths}}
\fancyhead[C]{{\fontfamily{phv}\selectfont $\cdots$ Colle n°3 - Léandre Courtoy $\cdots$}}
\fancyhead[R]{{\fontfamily{phv}\selectfont Page \thepage /\pageref{LastPage}}}
{\fontfamily{phv}\selectfont
\begin{center}
\Huge{Colle n°3} %Titre à changer
\end{center}
%Alors c'est techniquement la première, mais c'est celle de semaine 3
\section{Automatismes}
\noindent \textit{cf}: \href{https://1drv.ms/b/c/7e66df76a526ba79/IQBagJNN3ihvT7w402za7VA-AcGb2SWqQ4ZRgtFisU0i2aU?e=mMaqkx}{\pg{Nombres complexes.pdf}}
\section{Questions de cours}
\Question{Déterminer le domaine de définition et la limite en $+\infty$ de $\left(\left(1+\frac{2}{t}\right)^t,-t\sin \left(\frac1{t-1}\right)\right)$.}{Chap 1. Ex A2/A4}\\
\But{On détermine le domaine de définition de chaque composante, puis on étudie les limites en $+\infty$ à l'aide d'équivalents et de développements limités.}\\
Notons : $f_1:t\mapsto\left(1+\dfrac{2}{t}\right)^t$ et $f_2:t\mapsto-t\sin \left(\dfrac1{t-1}\right)$.\\

\begin{multicols}{2}
\noindent \textbf{Domaine de définition :}\\
Commençons par déterminer $\mathcal{D}_{f_2}$ :\\
Soit $t\in\mathbb{R}$,
\begin{align*}
f_2 \text{ est définie en } t &\Leftrightarrow \frac{1}{t-1} \text{ existe}\\
&\Leftrightarrow t-1\neq 0\\
&\Leftrightarrow t\neq 1
\end{align*}
Donc $\mathcal{D}_{f_2}=\mathbb{R}\setminus\{1\}$.

\columnbreak

Déterminons $\mathcal{D}_{f_1}$ :\\
Soit $t\in\mathbb{R}$,\\
$f_1(t)=\left(1+\dfrac{2}{t}\right)^t=\exp\left(t\ln\left(1+\frac{2}{t}\right)\right)$
\begin{align*}
f_1 \text{ est définie en } t &\Leftrightarrow \left\{\begin{matrix}
1+\frac{2}{t} \text{ existe} \\
\ln\left(1+\frac{2}{t}\right) \text{ existe}
\end{matrix}\right.\\
&\Leftrightarrow \left\{\begin{matrix}
t\neq 0\\
1+\frac{2}{t}>0
\end{matrix}\right.\\
&\Leftrightarrow\left\{\begin{matrix}
t\neq 0\\
\frac{t+2}{t}>0
\end{matrix}\right.
\end{align*}
D'après le tableau de signe de $\frac{t+2}{t}$ :
\begin{center}
\includegraphics[scale=0.5]{Tableau de signe.pdf} 
\end{center}
$\mathcal{D}_{f_1}=]-\infty,-2[\,\cup\,]0,+\infty[$.
\end{multicols}

\noindent Ainsi, $\mathcal{D}_f = \mathcal{D}_{f_1} \cap \mathcal{D}_{f_2} = ]-\infty,-2[\,\cup\,]0,1[\,\cup\,]1,+\infty[$.

\vspace{0.5cm}

\begin{multicols}{2}
\noindent\textbf{Limites en $+\infty$ :}\\
• Limite de $f_2$ :\\
\textcolor{myBlue}{Forme indéterminée "$\infty \times 0$"}\\
Comme $\displaystyle\lim_{t\to+\infty} \frac{1}{t-1} = 0$, on a :
$$\sin\left(\frac{1}{t-1}\right) \underset{t\to +\infty}{\sim} \frac{1}{t-1} \underset{t\to +\infty}{\sim} \frac{1}{t}$$
D'où :
$$f_2(t) = -t\sin\left(\frac{1}{t-1}\right) \underset{t\to +\infty}{\sim} -t \times \frac{1}{t} = -1$$
Donc : $\displaystyle\lim_{t\to +\infty} f_2(t) = -1$.

\columnbreak

\noindent • Limite de $f_1$ :\\
Posons $x = \frac{2}{t} \xrightarrow[t\to+\infty]{} 0$.\\
Pour $t > 0$, on peut effectuer un DL à l'ordre 2:
\begin{align*}
\ln\left(1+\frac{2}{t}\right) \underset{t\to+\infty}{=}& \frac{2}{t} - \frac{2}{t^2} + o\left(\frac{1}{t^2}\right)\\
 \underset{t\to+\infty}{=}& 2 - \frac{2}{t} + o\left(\frac{1}{t}\right)
\end{align*}
Comme $\displaystyle\lim_{t\to+\infty} \left(2 - \frac{2}{t} + o\left(\frac{1}{t}\right)\right) = 2$, par continuité de la fonction exponentielle en $2$ :
$$\lim_{t\to +\infty} f_1(t) = e^2$$
\end{multicols}

\noindent $f_1$ et $f_2$ admettent des limites finies en $+\infty$, donc :
$$\boxed{\lim_{t\to +\infty} f(t) = (e^2, -1)}\pagebreak$$
\Question{Si $f,g\in\mathcal{C}^1(I,\R^2)$, montrer que $\psi:t\mapsto\det(f(t),g(t))$ est $\mathcal{C}^1$ sur $I$ et déterminer l'expression de sa dérivée.\phantom{blablablabldsfdfsdfsdfabal} \phantom{ablablablab blabla}}{Chap 1. Prop B8}\\
\But{En gros on dérive le produit d'une fonction qui est la composante de 4 sous fonctions (venant de 2 couples) pour en retrouver l'expression de 2 déterminants qu'on sort par magie.. merci les maths !!}
\begin{demomethode}{}{\begin{itemize}
\item On prend un $t$ dans $\R$
\item On pose $f(t)$ et $g(t)$ tout deux des couples de 2 sous fonctions
\iitem Ces 2 fonct$^\circ$ sont $\mathcal{C}^1(\R,\R)$
\textit{Sinon on va pas pouvoir utiliser leurs dérivées et ce serait idiot c'est ce qu'on cherche..}
\item On pose $\psi$ le $\det$ de $f$ et $g$
\item On exprime le $\det$ (règle $\gamma$)
\item On a des produits de fonctions $\mathcal{C}^1$, donc $\psi$ est continue et dérivable sur $\R$
\item Dériver $\psi$
\item Exprimer les $\det$ de $f',g$ et $g',f$
\item Reconnaitre de l'expression $\psi'$ nos $\det$
\item D'où $\psi'=\det(f',g)+\det(f,g')$
\end{itemize}}
\xBox{Soit $t\in\R$}{
	\xBox{Posons $f(t)=(f_1(t),f_2(t)),g(t)=(g_1(t),g_2(t))$}{
	On note $f_1,f_2,g_1,g_2$ des fonctions $\mathcal{C}^1$ sur $\R$
		\xBox{Posons $\psi:t\mapsto \det(f(t),g(t))$}{
		$$\psi(t)=\begin{vmatrix}f_1(t) & g_1(t) \\f_2(t) & g_2(t)\end{vmatrix}=f_1(t)g_2(t)-f_2(t)g_1(t)$$
		$\psi$ est de classe $\mathcal{C}^1$ sur $\R$ par \textbf{produit} et:
		$$\psi'(t)=f_1'(t)g_2(t)+f_1g_2'(t)-\left[f_2'\left(t\right)g_1(t)+f_2(t)g_1'(t)\right]$$
		Or, $$\det(f'(t),g(t))=\begin{vmatrix}
		f_1'(t) & g_1(t)\\ f_2'(t) & g_2(t)\end{vmatrix}=f_1'(t)g_2(t)-f_2'(t)g_1(t)$$
		$$\det(f(t),g'(t))=\begin{vmatrix}
		f_1(t) & g_1'(t)\\ f_2(t) & g_2'(t)\end{vmatrix}=f_1(t)g_2'(t)-f_2(t)g_1'(t)$$
		On reconnait alors:
		\begin{align*}
		\psi'(t)=&f_1'(t)g_2(t)+f_1g_2'(t)-\left[f_2'\left(t\right)g_1(t)+f_2(t)g_1'(t)\right]\\
		=& f_1'(t)g_2(t)+f_1g_2'(t)-f_2'\left(t\right)g_1(t)-f_2(t)g_1'(t)\\
		=& \underbrace{f_1'(t)g_2(t)-f_2'\left(t\right)g_1(t)}_{\det(f'(t),g(t))}+\underbrace{f_1g_2'(t)-f_2(t)g_1'(t)}_{\det(f(t),g'(t)}
		\end{align*}
Ainsi,
$$\boxed{\psi'(t)=\det(f'(t),g(t))+\det(f(t),g'(t))}$$		
		}{et ce en posant $\psi:t\mapsto \det(f(t),g(t))$}	
	}{tout en posant $f(t)=(f_1(t),f_2(t)),g(t)=(g_1(t),g_2(t))$}
}{et ce pour tout $t\in\R$}
\end{demomethode}\pagebreak
\Question{On pose $\Gamma$ la courbe paramétrée par
$x(t)=\ln(2+t)+\dfrac{1}{t}\quad\textrm{et}\quad y(t)=t+\dfrac1t$.\\
Tracer l'allure de $\Gamma$ au voisinage de $M(1)$ ou $M(-1)$ \emph{(au choix de l'interrogateur)}.}{Chap 1. Ex C11}\\
\But{On fait des DL tout en pensant à bien factoriser afin d'observer l'allure locale}. \hfill Rappel: \href{https://1drv.ms/b/c/7e66df76a526ba79/IQDGI4ukOeMMRblxG_1sm0oSAflf_rlHMEyRKOFjhgNI8VM?e=RlXmOA}{\pg{Développements limités.pdf}}
\begin{demomethode}{Cas de $M(1)$}{
\begin{itemize}
\item On se place en $t=1+h$
\iitem ça nous permet de satisfaire le $\dfrac{1}{t}$
\item On fait nos DL
\iitem On pensera à factoriser ce qu'il y a dans notre $\ln$
\item On a $p=1$ et $q=2$, on reconnait donc un point ordinaire
\item On trace l'allure locale (parabole)\\\\
{\footnotesize Faire de même avec $M(-1)$ juste ce sera plus $t=1+h$ mais $t=-1+h$ et on aura un point de rebroussement de 1ère espèce}
\end{itemize}
}
Effectuons un DL à l'ordre 3 en $1$ de $x$ et $y$.\\
Soit $h$ au voisinage de $0$.
\begin{align*}
x(1+h)&\xeq{h\to 0} \ln(3+h)+\dfrac{1}{1+h}\\
&= \ln\left(3\left(1+\dfrac{h}{3}\right)\right)+\dfrac{1}{1+h}\\
&=\ln(3)+\dfrac{h}{3}-\dfrac{h^2}{3^2\cdot 2}+\dfrac{h^3}{3^3\cdot 3}+o(\dfrac{h^3}{3})+1-h+h^2-h^3+o\left(h^3\right)\\
x(1+h) &= \ln(3) + 1 - \frac{2}{3}h + \frac{17}{18}h^2 - \frac{80}{81}h^3 + o(h^3)
\end{align*}
\begin{align*}
y(1+h)&\xeq{h\to0} 1+h+\dfrac{1}{1+h}\\
&= 1+h+1-h+h^2-h^3+o\left(h^3\right)\\
&=2+h^2-h^3+o\left(h^3\right)
\end{align*}
Alors,\\
$M(1)$ a pour coordonnées $(1+\ln(3),2)$\\
On est dans le cas $p=1$ et $q=2$ et:
$$f^{(p)}(1)=\left(-\dfrac{2}{3},0\right) \quad \text{et} \quad f^{(q)}(1)=\left(\dfrac{17}{9},2\right)$$
Ainsi $M(1)$ est un \textbf{point ordinaire} et on a l'allure local de $\Gamma$ (la courbe)
\end{demomethode}
\begin{center}\vspace*{-30pt}
\includegraphics[scale=0.8]{Allure_locale_de_Gamma.pdf}\vspace*{-10pt}
\includegraphics[scale=0.5]{Allures_locales.pdf} 
\end{center}
\end{document}